\documentclass[12pt]{article}

% === CÀI ĐẶT TRANG ===
\usepackage[margin=2cm]{geometry}
\usepackage[utf8]{inputenc}
\usepackage[T5]{fontenc}
\usepackage[vietnamese]{babel}

% === TOÁN HỌC & HÌNH VẼ ===
\usepackage{amsmath, amssymb, amsthm}
\usepackage{graphicx}
\usepackage{float}
\usepackage{tikz}
\usepackage{pgfplots}
\pgfplotsset{compat=1.18}

% === ĐỊNH DẠNG ===
\usepackage{enumitem}
\usepackage{multicol}
\usepackage{fancyhdr}
\usepackage[hidelinks]{hyperref}

% Thiết lập khoảng cách cho list để văn bản thoáng hơn
\setlist[enumerate]{itemsep=4pt, topsep=4pt}

% === HEADER/FOOTER ===
\pagestyle{fancy}
\fancyhf{}
\fancyhead[L]{\small Đề cương ôn tập KTTX 1}
\fancyhead[R]{\small Lớp 2026 - L01}
\fancyfoot[C]{Trang \thepage}
\renewcommand{\headrulewidth}{0.4pt}

\begin{document}
	
	% ================================================
	% TRANG BÌA
	% ================================================
	\begin{titlepage}
		\centering
		\vspace*{4cm}
		
		{\Huge \textbf{ĐỀ CƯƠNG ÔN TẬP}}\\[1cm]
		{\LARGE \textbf{KIỂM TRA THƯỜNG XUYÊN 1}}\\[1.5cm]
		
		{\Large \textbf{Lớp: 2026 - L01}}\\[0.5cm]
		
		{\large \textit{(Lưu hành nội bộ)}}
		
		\vfill
	\end{titlepage}
	
	% ================================================
	% ĐỀ SỐ 1
	% ================================================
	\newpage % Bắt buộc sang trang mới
	\begin{center}
		{\LARGE \textbf{ĐỀ 1}}\\[4pt]
	\end{center}
	\vspace{4mm}
	
	% --- PHẦN 1: TRẮC NGHIỆM ---
	\noindent\textbf{PHẦN I. Câu trắc nghiệm nhiều phương án lựa chọn.} Thí sinh trả lời từ câu 1 đến câu 12. Mỗi câu hỏi thí sinh chỉ chọn một phương án.
	\vspace{2mm}
	% --- Bắt đầu từ dưới dòng \noindent\textbf{PHẦN I...} ---
	
	\begin{enumerate}[label=\textbf{Câu \arabic*.}, leftmargin=*]
		
		% CÂU 1 (Chế từ ảnh 1 - Tiệm cận đứng)
		\item Hàm số nào dưới đây có đồ thị \textbf{không} có tiệm cận đứng?
		\begin{multicols}{4}
			\begin{enumerate}[label=\Alph*.]
				\item $y = \dfrac{x+2}{x-2}$.
				\item $y = \dfrac{x-1}{x^2+2}$.
				\item $y = \dfrac{x^2-3x}{x+1}$.
				\item $y = \dfrac{x}{x^2-4}$.
			\end{enumerate}
		\end{multicols}
		
		% CÂU 2 (Chế từ ảnh 2 - Bảng biến thiên)
		\item Cho hàm số $y = f(x)$ có bảng biến thiên như sau:
		\begin{center}
			\begin{tikzpicture}[xscale=1.2, yscale=0.8]
				% Khung bảng
				\draw (0,0) rectangle (8,3);
				\draw (0,1) -- (8,1);
				\draw (0,2) -- (8,2);
				\draw (1,0) -- (1,3);
				
				% Dòng x
				\node at (0.5, 2.5) {$x$};
				\node at (1.5, 2.5) {$-\infty$};
				\node at (3.5, 2.5) {$-1$};
				\node at (5.5, 2.5) {$2$};
				\node at (7.5, 2.5) {$+\infty$};
				
				% Dòng y'
				\node at (0.5, 1.5) {$y'$};
				\node at (2.5, 1.5) {$-$};
				\node at (3.5, 1.5) {$0$};
				\node at (4.5, 1.5) {$+$};
				\node at (5.5, 1.5) {$0$};
				\node at (6.5, 1.5) {$-$};
				
				% Dòng y
				\node at (0.5, 0.5) {$y$};
				\node at (1.5, 0.7) {$+\infty$};
				\node at (3.5, 0.2) {$-2$};
				\node at (5.5, 0.7) {$3$};
				\node at (7.5, 0.2) {$-\infty$};
				
				% Mũi tên
				\draw[->, thick, blue] (1.8, 0.6) -- (3.2, 0.3);
				\draw[->, thick, blue] (3.8, 0.3) -- (5.2, 0.6);
				\draw[->, thick, blue] (5.8, 0.6) -- (7.2, 0.3);
			\end{tikzpicture}
		\end{center}
		Hàm số đã cho đồng biến trên khoảng nào dưới đây?
		\begin{multicols}{4}
			\begin{enumerate}[label=\Alph*.]
				\item $(0; 1)$.
				\item $(2; 5)$.
				\item $(-2; 0)$.
				\item $(-\infty; -1)$.
			\end{enumerate}
		\end{multicols}
		
		% CÂU 3 (Chế từ ảnh 3 - Nhận diện đồ thị)
		\item Hàm số nào dưới đây có đồ thị như đường cong trong hình bên?
		
		\begin{center}
			\begin{tikzpicture}
				\begin{axis}[
					axis lines = middle,
					xlabel = {$x$}, ylabel = {$y$},
					ymin = -2, ymax = 4,
					xmin = -1.5, xmax = 3.5,
					xtick = {2}, ytick = {-1, 3},
					width=8cm, height=6cm,
					samples=100,
					axis line style={->}
					]
					% Đồ thị y = -x^3 + 3x^2 - 1
					\addplot [domain=-1.2:3.2, thick] {-x^3 + 3*x^2 - 1};
					% Nét đứt gióng tọa độ
					\addplot [dashed] coordinates {(2,0) (2,3) (0,3)};
					% Gốc tọa độ
					\node[below left] at (axis cs:0,0) {$O$};
				\end{axis}
			\end{tikzpicture}
		\end{center}
		
		\begin{center}
			\begin{tabular}{p{0.47\textwidth} p{0.47\textwidth}}
				A. $y=-x^3+3x^2-1$. &
				B. $y=x^3-3x^2-1$. \\[2mm]
				C. $y=\dfrac{-x+2}{x+1}$. &
				D. $y=-x^3-3x^2-1$.
			\end{tabular}
		\end{center}
		
		% CÂU 4 (Chế từ ảnh 4 - Giao điểm tiệm cận)
		\item Tọa độ tâm đối xứng $I$ của đồ thị hàm số $y = \dfrac{x^2+2x-2}{x+1}$ là:
		\begin{multicols}{4}
			\begin{enumerate}[label=\Alph*.]
				\item $I(-1; 0)$.
				\item $I(1; 2)$.
				\item $I(-1; 1)$.
				\item $I(0; -1)$.
			\end{enumerate}
		\end{multicols}
		
		% CÂU 5 (Chế từ ảnh 1 - Điểm cực trị của hàm bậc 3)
		\item Đồ thị hàm số $y = x^3 + 3x^2 - 9x - 2$ có tọa độ điểm cực tiểu là:
		\begin{multicols}{4}
			\begin{enumerate}[label=\Alph*.]
				\item $(1; -7)$.
				\item $(-3; 25)$.
				\item $(1; 7)$.
				\item $(-3; -7)$.
			\end{enumerate}
		\end{multicols}
		
		% CÂU 6 (Chế từ ảnh 2 - Đọc đồ thị f'(x) tìm khoảng đơn điệu)
		\item Cho hàm số $y = f(x)$ liên tục trên $\mathbb{R}$. Hàm số đạo hàm $y = f'(x)$ có đồ thị như hình vẽ bên dưới.
		\begin{center}
			\begin{tikzpicture}
				\begin{axis}[
					axis lines = middle,
					xlabel = {$x$}, ylabel = {$y$},
					ymin = -3, ymax = 3,
					xmin = -2.5, xmax = 4.5,
					xtick = {-1, 3}, ytick = \empty,
					xticklabel style = {xshift=0mm, yshift=-2mm, red},
					width=8cm, height=5cm,
					samples=100,
					axis line style={->}
					]
					% Đồ thị f'(x) (parabol cắt trục hoành tại -1 và 3)
					\addplot [domain=-1.8:3.8, very thick, blue] {0.5*(x+1)*(x-3)} node[pos=0.9, left] {$y = f'(x)$};
					
					% Điểm cắt trục hoành
					\addplot[mark=*, mark size=1.5pt, red] coordinates {(-1,0)};
					\addplot[mark=*, mark size=1.5pt, red] coordinates {(3,0)};
					
					% Gốc tọa độ
					\node[below left] at (axis cs:0,0) {$O$};
				\end{axis}
			\end{tikzpicture}
		\end{center}
		Hàm số $y = f(x)$ nghịch biến trên khoảng nào dưới đây?
		\begin{multicols}{4}
			\begin{enumerate}[label=\Alph*.]
				\item $(-1; 3)$.
				\item $(-\infty; -1)$.
				\item $(3; +\infty)$.
				\item $(0; 4)$.
			\end{enumerate}
		\end{multicols}
		
		% CÂU 7 (Chế từ ảnh 3 - Min/Max trên đoạn)
		\item Tìm giá trị nhỏ nhất của hàm số $y = \dfrac{3x-1}{x+1}$ trên đoạn $[0; 2]$.
		\begin{multicols}{4}
			\begin{enumerate}[label=\Alph*.]
				\item $-1$.
				\item $\dfrac{5}{3}$.
				\item $2$.
				\item $0$.
			\end{enumerate}
		\end{multicols}
		
		% CÂU 8 (Yêu cầu thêm - Khai triển đạo hàm hàm đa thức)
		\item Cho hàm số $y = \dfrac{1}{3}x^3 - 2x^2 + 3x - 5$. Đạo hàm $y'$ của hàm số đã cho là:
		
		\begin{center}
			\begin{tabular}{p{0.45\textwidth} p{0.45\textwidth}}
				A. $y' = x^2 - 4x + 3$. &
				B. $y' = x^2 - 2x + 3$. \\[2mm]
				C. $y' = 3x^2 - 4x + 3$. &
				D. $y' = \dfrac{1}{3}x^2 - 4x + 3$.
			\end{tabular}
		\end{center}
		
		
		% CÂU 9 (Chế từ ảnh 1 - Max/Min trên đoạn từ đồ thị)
		\item Cho hàm số $y = f(x)$ liên tục trên đoạn $[1; 4]$ và có đồ thị như hình vẽ bên dưới.
		\begin{center}
			\begin{tikzpicture}[scale=0.9]
				% Lưới mờ (tùy chọn, giúp giống đồ thị gốc)
				\draw[step=1cm, gray!20, very thin, dashed] (0,-2) grid (5,4);
				
				% Hệ trục tọa độ
				\draw[->] (-0.5, 0) -- (5.5, 0) node[right] {$x$};
				\draw[->] (0, -2.5) -- (0, 4.5) node[above] {$y$};
				\node[below left] at (0,0) {$O$};
				
				% Các điểm trên trục
				\foreach \x in {1, 2, 4} \draw (\x, 0.1) -- (\x, -0.1) node[below] {$\x$};
				\foreach \y in {-1, 2, 3} \draw (0.1, \y) -- (-0.1, \y) node[left] {$\y$};
				
				% Đường gióng
				\draw[dashed, gray] (1,0) -- (1,2) -- (0,2);
				\draw[dashed, gray] (2,0) -- (2,3) -- (0,3);
				\draw[dashed, gray] (4,0) -- (4,-1) -- (0,-1);
				
				% Đồ thị parabol ngược y = -x^2 + 4x - 1 trên [1;4]
				\draw[very thick, blue, domain=1:4, samples=100] plot (\x, {-\x*\x + 4*\x - 1});
				
				% Điểm nhấn (đỏ)
				\fill[red] (1,2) circle (2pt);
				\fill[red] (2,3) circle (2pt);
				\fill[red] (4,-1) circle (2pt);
			\end{tikzpicture}
		\end{center}
		Gọi $M, m$ lần lượt là giá trị lớn nhất và giá trị nhỏ nhất của hàm số trên đoạn $[1; 4]$. Tính giá trị của biểu thức $T = M + m$.
		\begin{multicols}{4}
			\begin{enumerate}[label=\Alph*.]
				\item $T = 2$.
				\item $T = 4$.
				\item $T = 1$.
				\item $T = 5$.
			\end{enumerate}
		\end{multicols}
		
		% CÂU 10 (Chế từ ảnh 2 - Số điểm cực trị từ phương trình f'(x))
		\item Cho hàm số $y = f(x)$ có đạo hàm $f'(x) = x^2(x-2)^3(x+1)^5(x-5)$ với mọi $x \in \mathbb{R}$. Số điểm cực trị của hàm số đã cho là:
		\begin{multicols}{4}
			\begin{enumerate}[label=\Alph*.]
				\item $4$.
				\item $2$.
				\item $3$.
				\item $1$.
			\end{enumerate}
		\end{multicols}
		
		% CÂU 11 (Chế từ ảnh 3 - Nhìn Bảng biến thiên tìm giá trị cực trị)
		\item Cho hàm số $y = f(x)$ liên tục trên $\mathbb{R}$ và có bảng biến thiên như sau:
		\begin{center}
			\begin{tikzpicture}[xscale=1.2, yscale=0.8]
				% Khung bảng
				\draw (0,0) rectangle (10,3);
				\draw (0,1) -- (10,1);
				\draw (0,2) -- (10,2);
				\draw (1,0) -- (1,3);
				
				% Dòng x
				\node at (0.5, 2.5) {$x$};
				\node at (1.5, 2.5) {$-\infty$};
				\node at (3.5, 2.5) {$-1$};
				\node at (5.5, 2.5) {$1$};
				\node at (7.5, 2.5) {$3$};
				\node at (9.5, 2.5) {$+\infty$};
				
				% Dòng y'
				\node at (0.5, 1.5) {$y'$};
				\node at (2.5, 1.5) {$+$};
				\node at (3.5, 1.5) {$0$};
				\node at (4.5, 1.5) {$-$};
				\node at (5.5, 1.5) {$0$};
				\node at (6.5, 1.5) {$+$};
				\node at (7.5, 1.5) {$0$};
				\node at (8.5, 1.5) {$-$};
				
				% Dòng y
				\node at (0.5, 0.5) {$y$};
				\node at (1.5, 0.2) {$-\infty$};
				\node at (3.5, 0.7) {$5$};
				\node at (5.5, 0.2) {$2$};
				\node at (7.5, 0.7) {$5$};
				\node at (9.5, 0.2) {$-\infty$};
				
				% Mũi tên
				\draw[->, thick, blue] (1.8, 0.3) -- (3.2, 0.6);
				\draw[->, thick, blue] (3.8, 0.6) -- (5.2, 0.3);
				\draw[->, thick, blue] (5.8, 0.3) -- (7.2, 0.6);
				\draw[->, thick, blue] (7.8, 0.6) -- (9.2, 0.3);
			\end{tikzpicture}
		\end{center}
		Giá trị cực tiểu của hàm số đã cho bằng bao nhiêu?
		\begin{multicols}{4}
			\begin{enumerate}[label=\Alph*.]
				\item $x = 1$.
				\item $y = 5$.
				\item $y = 2$.
				\item $x = -1$.
			\end{enumerate}
		\end{multicols}
		
		% CÂU 12 (Chế từ ảnh 4 - Khoảng đơn điệu của hàm phân thức bậc 2/bậc 1)
		\item Hỏi hàm số $y = \dfrac{x^2 + 2x + 2}{x + 1}$ nghịch biến trên các khoảng nào dưới đây?
		
		\begin{center}
			\begin{tabular}{p{0.47\textwidth} p{0.47\textwidth}}
				A. $(-2; 0)$. &
				B. $(-\infty; -2)$ và $(0; +\infty)$. \\[2mm]
				C. $(-2; -1)$ và $(-1; 0)$. &
				D. $(-1; 0)$.
			\end{tabular}
		\end{center}
	\end{enumerate}
	\vspace{4mm}
	% --- PHẦN 2: ĐÚNG / SAI ---
	\noindent\textbf{PHẦN II. Câu trắc nghiệm đúng sai.} Thí sinh trả lời từ câu 1 đến câu 2. Trong mỗi ý a), b), c), d) ở mỗi câu, thí sinh chọn đúng hoặc sai.
	\vspace{2mm}
	
	\begin{enumerate}[label=\textbf{Câu \arabic*.}, leftmargin=*]
			
		% CÂU 1 ĐÚNG/SAI (Chế từ ảnh - Hàm phân thức bậc 2/bậc 1)
		\item Cho hàm số $y = \dfrac{ax^2+bx+c}{mx+n}$ ($a \neq 0, m \neq 0$) với $a, b, c, m, n$ là các số thực và có đồ thị $(C)$ như hình vẽ dưới đây.
		
		\begin{center}
			\begin{tikzpicture}
				\begin{axis}[
					axis lines = middle,
					xlabel = {$x$}, ylabel = {$y$},
					xmin = -2.5, xmax = 4.5,
					ymin = -3.5, ymax = 4.5,
					xtick = {1, 2}, ytick = {-2, 2},
					ticklabel style = {fill=white, inner sep=1pt},
					width=9cm, height=8cm,
					samples=100,
					axis line style={->}
					]
					
					% Tiệm cận xiên y = x - 1
					\addplot[domain=-2.5:4.5, red, thick, dashed] {x-1} node[pos=0.85, below right, black] {$y = x - 1$};
					
					% Tiệm cận đứng x = 1
					\draw[red, thick, dashed] (1, -3.5) -- (1, 4.5);
					
					% Đồ thị nhánh trái (x < 1)
					\addplot[domain=-2.5:0.75, blue, thick] {x - 1 + 1/(x-1)};
					
					% Đồ thị nhánh phải (x > 1)
					\addplot[domain=1.25:4.5, blue, thick] {x - 1 + 1/(x-1)};
					
					% Gióng tọa độ điểm cực tiểu (2, 2)
					\draw[gray, dashed] (2, 0) -- (2, 2) -- (0, 2);
					\fill (2, 2) circle (1.5pt);
					
					% Điểm cực đại (0, -2) nằm trên trục Oy
					\fill (0, -2) circle (1.5pt);
					
					% Gốc O
					\node[below left] at (0,0) {$O$};
				\end{axis}
			\end{tikzpicture}
		\end{center}
		
		Các mệnh đề sau đúng hay sai?
		\begin{enumerate}[label=\alph*)]
			\item Hàm số $y = f(x)$ đồng biến trên các khoảng $(-\infty; 0)$ và $(2; +\infty)$.
			\item Đồ thị hàm số có tiệm cận đứng là $x = 1$ và tiệm cận xiên là $y = x - 1$.
			\item Đồ thị hàm số có điểm cực đại là $(2; 2)$ và điểm cực tiểu là $(0; -2)$.
			\item Trên khoảng $(1; +\infty)$, giá trị nhỏ nhất của hàm số $y = f(x)$ bằng $2$.
		\end{enumerate}
		
		% CÂU 2 ĐÚNG/SAI (Chế từ ảnh - Ứng dụng đạo hàm trong chuyển động)
		\item Một vật chuyển động thẳng có phương trình xác định bởi:
		\[ s(t) = t^3 - 6t^2 + 9t + 5 \quad (t \ge 0) \]
		trong đó $t$ tính bằng giây và $s(t)$ tính bằng mét. Xét tính đúng hay sai của các mệnh đề sau:
		\begin{enumerate}[label=\alph*)]
			\item Vận tốc tức thời của vật bằng $0$ tại các thời điểm $t = 1$ giây và $t = 3$ giây.
			\item Phương trình biểu diễn gia tốc tức thời của vật là $a(t) = 6t - 12 \text{ (m/s}^2\text{)}$.
			\item Vận tốc của vật giảm khi $t > 2$ giây.
			\item Vận tốc của vật giảm trong suốt khoảng thời gian từ $1$ giây đến $3$ giây ($t \in (1; 3)$).
		\end{enumerate}
		
		
	\end{enumerate}
	
	\vspace{4mm}
	% --- PHẦN 3: TRẢ LỜI NGẮN ---
	\noindent\textbf{PHẦN III. Câu trắc nghiệm trả lời ngắn.} Thí sinh trả lời từ câu 1 đến câu 4.
	\vspace{2mm}
	
	% --- Bắt đầu từ dưới dòng \noindent\textbf{PHẦN III...} ---
	
	\begin{enumerate}[label=\textbf{Câu \arabic*.}, leftmargin=*]
		
		% CÂU 1 TRẢ LỜI NGẮN (Chế từ ảnh 1 - Khoảng đồng biến)
		\item Biết hàm số $y = \dfrac{x^2 - 4x + 8}{x - 2}$ đồng biến trên khoảng $(a; +\infty)$. Tính giá trị của $a$.
		
		% CÂU 2 TRẢ LỜI NGẮN (Chế từ ảnh 2 - Max hàm phân thức trên đoạn)
		\item Tìm giá trị lớn nhất của hàm số $y = \dfrac{3x - 2}{x + 1}$ trên đoạn $[0; 3]$ (viết kết quả dưới dạng số thập phân).
		
		% CÂU 3 TRẢ LỜI NGẮN (Chế từ ảnh 3 - Bài toán thực tế G(x))
		\item Độ giảm huyết áp của một bệnh nhân được cho bởi công thức $G(x) = 0.05x^2(24 - x)$ trong đó $x$ là liều lượng thuốc được tiêm cho bệnh nhân (tính bằng miligam, $0 < x < 24$). Biết liều lượng thuốc cần tiêm nằm trong khoảng $(a;b)$ thì độ giảm huyết áp của bệnh nhân tăng dần. Tính giá trị của tổng $a + b$.
		
		% CÂU 4 TRẢ LỜI NGẮN (Chế từ ảnh 4 - Độ dài đoạn nối 2 điểm cực trị)
		\item Cho hàm số $y = \dfrac{x^2 - 6x + 13}{x - 3}$ có đồ thị là đường cong $(C)$. Tính độ dài đoạn thẳng nối hai điểm cực trị của đồ thị $(C)$ (làm tròn kết quả đến hàng phần mười).
		
	\end{enumerate}
	
	\vspace{4mm}
	% --- PHẦN 4: TỰ LUẬN ---
	\noindent\textbf{PHẦN IV. Tự luận.} Thí sinh trình bày bài giải chi tiết từ bài 1 đến bài 3.
	\vspace{2mm}
	
	% --- Bắt đầu từ dưới dòng \noindent\textbf{PHẦN IV...} ---
	
	\begin{enumerate}[label=\textbf{Bài \arabic*.}, leftmargin=*]
		
		% BÀI 1 (Chế từ ảnh 1 - Bài toán Max/Min thực tế)
		\item Vận tốc của một thiết bị bay không người lái (drone) từ lúc cất cánh ($t=0$ s) cho đến khi hạ cánh được mô hình hóa bởi phương trình $v(t) = 0.002t^3 - 0.24t^2 + 20$ (trong đó $v$ tính bằng m/s, $t$ tính bằng giây). Hỏi tại thời điểm $t$ bằng bao nhiêu giây kể từ lúc cất cánh thì gia tốc của thiết bị bay đạt giá trị nhỏ nhất?
		
		\vspace{2mm}
		% BÀI 2 (Chế từ ảnh 2 - Yêu cầu dùng Heron tính diện tích)
		\item Cho hàm số $y = \dfrac{x^2 + 2x + 2}{x + 1}$ có đồ thị là đường cong $(C)$. Gọi $A, B$ lần lượt là hai điểm cực trị của đồ thị $(C)$ và điểm $C(0; -2)$. Tính diện tích tam giác ABC.
		
		\vspace{2mm}
		% BÀI 3 (Chế từ ảnh 3 - Khảo sát hàm số không vẽ hình)
		\item Khảo sát sự biến thiên (không vẽ bảng biến thiên), đồng thời xác định tâm đối xứng của hàm số $y = \dfrac{x^2 - 3x + 3}{x - 2}$. \\
		
	\end{enumerate}
	
	
	% ================================================
	% ĐỀ SỐ 2
	% ================================================
	\newpage % Bắt buộc sang trang mới
	\begin{center}
		{\LARGE \textbf{ĐỀ SỐ 2}}\\[4pt]
	\end{center}
	\vspace{4mm}
	
	% ================================================
	% ĐỀ SỐ 2 - PHẦN I: TRẮC NGHIỆM
	% ================================================
	\noindent\textbf{PHẦN I. Câu trắc nghiệm nhiều phương án lựa chọn.} Thí sinh trả lời từ câu 1 đến câu 12. Mỗi câu hỏi thí sinh chỉ chọn một phương án.
	\vspace{2mm}
	
	\begin{enumerate}[label=\textbf{Câu \arabic*.}, leftmargin=*]
		
		% CÂU 1: ĐẠO HÀM (Mức độ Nhận biết/Thông hiểu)
		\item Cho hàm số $y = \sqrt{x^2 - 2x + 3}$. Đạo hàm $y'$ của hàm số đã cho là:
		\begin{multicols}{2}
			\begin{enumerate}[label=\Alph*.]
				\item $y' = \dfrac{2x-2}{\sqrt{x^2-2x+3}}$.
				\item $y' = \dfrac{x-1}{\sqrt{x^2-2x+3}}$.
				\item $y' = \dfrac{1}{2\sqrt{x^2-2x+3}}$.
				\item $y' = (2x-2)\sqrt{x^2-2x+3}$.
			\end{enumerate}
		\end{multicols}
		
		% CÂU 2: ĐƠN ĐIỆU - Nhận biết đọc đồ thị
		\item Cho hàm số $y = f(x)$ liên tục trên $\mathbb{R}$ và có đồ thị như hình vẽ bên dưới.
		\begin{center}
			\begin{tikzpicture}[scale=0.8]
				% Hệ trục tọa độ
				\draw[->] (-2.5,0) -- (2.5,0) node[right] {$x$};
				\draw[->] (0,-3) -- (0,3) node[above] {$y$};
				\node[below left] at (0,0) {$O$};
				% Đồ thị y = x^3 - 3x
				\draw[domain=-1.85:1.85, smooth, variable=\x, blue, very thick] plot ({\x}, {\x*\x*\x - 3*\x});
				% Đường gióng cực trị
				\draw[dashed, gray] (-1,0) -- (-1,2) -- (0,2);
				\draw[dashed, gray] (1,0) -- (1,-2) -- (0,-2);
				\node[above] at (-1,0) {$-1$};
				\node[below] at (1,0) {$1$};
				\node[left] at (0,2) {$2$};
				\node[right] at (0,-2) {$-2$};
			\end{tikzpicture}
		\end{center}
		Hàm số đã cho nghịch biến trên khoảng nào dưới đây?
		\begin{multicols}{4}
			\begin{enumerate}[label=\Alph*.]
				\item $(-\infty; -1)$.
				\item $(1; +\infty)$.
				\item $(-\infty; 1)$.
				\item $(-1; 1)$.
			\end{enumerate}
		\end{multicols}
		
		% CÂU 3: ĐƠN ĐIỆU - Thông hiểu hàm phân thức bậc 1/bậc 1
		\item Cho hàm số $y = \dfrac{2x - 3}{x + 1}$. Phát biểu nào sau đây đúng?
		\begin{enumerate}[label=\Alph*.]
			\item Hàm số đồng biến trên các khoảng $(-\infty; -1)$ và $(-1; +\infty)$.
			\item Hàm số đồng biến trên tập số thực $\mathbb{R} \setminus \{-1\}$.
			\item Hàm số nghịch biến trên các khoảng $(-\infty; -1)$ và $(-1; +\infty)$.
			\item Hàm số đồng biến trên tập hợp $(-\infty; -1) \cup (-1; +\infty)$.
		\end{enumerate}
		
		% CÂU 4
		\item Hàm số $y = \dfrac{x^2 - 2x + 2}{x - 1}$ nghịch biến trên các khoảng nào dưới đây?
		
		\begin{center}
			\begin{tabular}{p{0.47\textwidth} p{0.47\textwidth}}
				A. $(-\infty; 0)$ và $(2; +\infty)$. &
				B. $(0; 2)$. \\[2mm]
				C. $(0; 1)$ và $(1; 2)$. &
				D. $(-\infty; 1)$ và $(1; +\infty)$.
			\end{tabular}
		\end{center}
		
		
		% CÂU 5: CỰC TRỊ - Nhận biết từ bảng biến thiên
		\item Cho hàm số $y = f(x)$ liên tục trên $\mathbb{R}$ và có bảng biến thiên dưới đây:
		\begin{center}
			\begin{tikzpicture}[xscale=1.1, yscale=0.8]
				\draw (0,0) rectangle (8,3);
				\draw (0,1) -- (8,1);
				\draw (0,2) -- (8,2);
				\draw (1,0) -- (1,3);
				
				\node at (0.5, 2.5) {$x$};
				\node at (1.5, 2.5) {$-\infty$};
				\node at (3.5, 2.5) {$-1$};
				\node at (5.5, 2.5) {$3$};
				\node at (7.5, 2.5) {$+\infty$};
				
				\node at (0.5, 1.5) {$y'$};
				\node at (2.5, 1.5) {$+$};
				\node at (3.5, 1.5) {$0$};
				\node at (4.5, 1.5) {$-$};
				\node at (5.5, 1.5) {$0$};
				\node at (6.5, 1.5) {$+$};
				
				\node at (0.5, 0.5) {$y$};
				\node at (1.5, 0.2) {$-\infty$};
				\node at (3.5, 0.7) {$4$};
				\node at (5.5, 0.2) {$-1$};
				\node at (7.5, 0.7) {$+\infty$};
				
				\draw[->, thick, blue] (1.8, 0.3) -- (3.2, 0.6);
				\draw[->, thick, blue] (3.8, 0.6) -- (5.2, 0.3);
				\draw[->, thick, blue] (5.8, 0.3) -- (7.2, 0.6);
			\end{tikzpicture}
		\end{center}
		Điểm cực tiểu của hàm số đã cho là:
		\begin{multicols}{4}
			\begin{enumerate}[label=\Alph*.]
				\item $x = -1$.
				\item $x = 3$.
				\item $y = -1$.
				\item $x = 0$.
			\end{enumerate}
		\end{multicols}
		
		% CÂU 6: CỰC TRỊ - Thông hiểu số cực trị từ f'(x) đa thức lũy thừa
		\item Cho hàm số $y = f(x)$ có đạo hàm $f'(x) = (x-1)^2(x+2)(x-3)^3(x+4)^4$ với mọi $x \in \mathbb{R}$. Hỏi hàm số $f(x)$ có bao nhiêu điểm cực trị?
		\begin{multicols}{4}
			\begin{enumerate}[label=\Alph*.]
				\item $2$.
				\item $4$.
				\item $3$.
				\item $1$.
			\end{enumerate}
		\end{multicols}
		
		% CÂU 7: CỰC TRỊ - Vận dụng hàm trị tuyệt đối y = |f(x)|
		\item Số điểm cực trị của hàm số $g(x) = |x^3 - 3x|$ là:
		\begin{multicols}{4}
			\begin{enumerate}[label=\Alph*.]
				\item $3$.
				\item $4$.
				\item $5$.
				\item $2$.
			\end{enumerate}
		\end{multicols}
		
		% CÂU 8: MAX MIN - Nhận biết từ đồ thị đoạn giới hạn
		\item Cho hàm số $y = f(x)$ liên tục trên đoạn $[-1; 3]$ và có đồ thị như hình vẽ bên dưới.
		\begin{center}
			\begin{tikzpicture}[scale=0.8]
				\draw[step=1cm, gray!20, very thin, dashed] (-2,-3) grid (4,5);
				\draw[->] (-2.2,0) -- (4.2,0) node[right] {$x$};
				\draw[->] (0,-2.5) -- (0,4.8) node[above] {$y$};
				\node[below left] at (0,0) {$O$};
				% Khảo sát vẽ đồ thị đi qua (-1,1), (1,4), (3,-2)
				\draw[very thick, blue, domain=-1:3, samples=100] plot (\x, {-0.5*\x*\x*\x + 1.5*\x + 3});
				
				\draw[dashed, gray] (-1,0) -- (-1,1) -- (0,1);
				\draw[dashed, gray] (1,0) -- (1,4) -- (0,4);
				\draw[dashed, gray] (3,0) -- (3,-2) -- (0,-2);
				
				\node[below] at (-1,0) {$-1$};
				\node[below] at (1,0) {$1$};
				\node[above] at (3,0) {$3$};
				\node[left] at (0,1) {$1$};
				\node[left] at (0,4) {$4$};
				\node[left] at (0,-2) {$-2$};
				
				\fill[red] (-1,1) circle (2pt);
				\fill[red] (1,4) circle (2pt);
				\fill[red] (3,-2) circle (2pt);
			\end{tikzpicture}
		\end{center}
		Gọi $M, m$ lần lượt là giá trị lớn nhất và giá trị nhỏ nhất của hàm số trên đoạn $[-1; 3]$. Tính giá trị của biểu thức $H = M - m$.
		\begin{multicols}{4}
			\begin{enumerate}[label=\Alph*.]
				\item $2$.
				\item $5$.
				\item $4$.
				\item $6$.
			\end{enumerate}
		\end{multicols}
		
		% CÂU 9: MAX MIN - Thông hiểu hàm bậc 1/bậc 1 trên đoạn
		\item Tìm giá trị lớn nhất của hàm số $y = \dfrac{3x - 1}{x - 2}$ trên đoạn $[3; 5]$.
		\begin{multicols}{4}
			\begin{enumerate}[label=\Alph*.]
				\item $8$.
				\item $\dfrac{14}{3}$.
				\item $5$.
				\item $14$.
			\end{enumerate}
		\end{multicols}
		
		% CÂU 10: TIỆM CẬN (Nhận biết/Thông hiểu công thức)
		\item Đường tiệm cận ngang của đồ thị hàm số $y = \dfrac{3 - 2x}{x + 1}$ là đường thẳng có phương trình:
		\begin{multicols}{4}
			\begin{enumerate}[label=\Alph*.]
				\item $y = 3$.
				\item $x = -1$.
				\item $y = -2$.
				\item $y = -\dfrac{3}{2}$.
			\end{enumerate}
		\end{multicols}
		
		% CÂU 11: TÂM ĐỐI XỨNG CỦA HÀM SỐ (Thông hiểu)
		\item Tìm tọa độ tâm đối xứng $I$ của đồ thị hàm số bậc nhất trên bậc nhất $y = \dfrac{x-3}{x+1}$.
		\begin{multicols}{4}
			\begin{enumerate}[label=\Alph*.]
				\item $I(1; -1)$.
				\item $I(-1; 1)$.
				\item $I(-1; -3)$.
				\item $I(1; 1)$.
			\end{enumerate}
		\end{multicols}
		
		% CÂU 12: NHẬN DIỆN HÀM SỐ (Thông hiểu đồ thị phân thức)
		\item Đường cong trong hình vẽ dưới đây là đồ thị của hàm số nào?
		\begin{center}
			\begin{tikzpicture}[scale=0.8]
				\draw[->] (-3,0) -- (4,0) node[right] {$x$};
				\draw[->] (0,-2) -- (0,5) node[above] {$y$};
				\node[below left] at (0,0) {$O$};
				
				% Các tiệm cận vẽ mờ dạng nét đứt
				\draw[red, dashed, thick] (1,-2) -- (1,5);
				\draw[red, dashed, thick] (-3,2) -- (4,2);
				\node[below, red] at (1.2,0) {$1$};
				\node[left, red] at (0,2.2) {$2$};
				
				% Nhánh bên trái tiệm cận đứng
				\draw[domain=-3:0.65, smooth, variable=\x, blue, very thick] plot (\x, {(2*\x - 1)/(\x - 1)});
				% Nhánh bên phải tiệm cận đứng
				\draw[domain=1.35:4, smooth, variable=\x, blue, very thick] plot (\x, {(2*\x - 1)/(\x - 1)});
				
				% Điểm mút cắt trục
				\fill (0,1) circle (2pt);
				\node[left] at (0,1) {$1$};
				\fill (2,3) circle (2pt);
				\draw[dashed, gray] (2,0) node[below] {$2$} -- (2,3) -- (0,3) node[left] {$3$};
			\end{tikzpicture}
		\end{center}
		\begin{multicols}{2}
			\begin{enumerate}[label=\Alph*.]
				\item $y = \dfrac{2x+1}{x-1}$.
				\item $y = \dfrac{x-1}{x-2}$.
				\item $y = \dfrac{2x-1}{x-1}$.
				\item $y = \dfrac{x+1}{x-1}$.
			\end{enumerate}
		\end{multicols}
		
	\end{enumerate}
	
	% ================================================
	% ĐỀ SỐ 2 - PHẦN II: ĐÚNG / SAI
	% ================================================
	\vspace{4mm}
	\noindent\textbf{PHẦN II. Câu trắc nghiệm đúng sai.} Thí sinh trả lời từ câu 1 đến câu 2. Trong mỗi ý a), b), c), d) ở mỗi câu, thí sinh chọn đúng hoặc sai.
	\vspace{2mm}
	
	\begin{enumerate}[label=\textbf{Câu \arabic*.}, leftmargin=*]
		
		% CÂU 1: KHẢO SÁT ĐỒ THỊ HÀM SỐ (Hàm đa thức bậc 3)
		\item Cho hàm số $y = f(x) = -x^3 + 3x + 2$ có đồ thị như hình vẽ bên dưới.
		\begin{center}
			\begin{tikzpicture}[scale=0.8]
				% Hệ trục tọa độ
				\draw[->] (-2.5,0) -- (3.5,0) node[right] {$x$};
				\draw[->] (0,-1) -- (0,5.5) node[above] {$y$};
				\node[below, left] at (0,0) {$O$};
				
				% Đồ thị hàm bậc 3: y = -x^3 + 3x + 2
				\draw[domain=-2.1:2.2, smooth, variable=\x, blue, very thick] plot ({\x}, {-\x*\x*\x + 3*\x + 2});
				
				% Đường nét đứt gióng tọa độ cực trị
				\draw[dashed, gray] (1,0) -- (1,4) -- (0,4);
				\fill (1,4) circle (2pt);
				\fill (-1,0) circle (2pt);
				
				% Điểm đặc biệt cắt Oy
				\fill (0,2) circle (1.5pt);
				
				% Các nhãn số liệu
				\node[below, red] at (-1,0) {$-1$};
				\node[below, red] at (1,0) {$1$};
				\node[below, red] at (2,0) {$2$};
				\node[left, red] at (0,4) {$4$};
				\node[left, red] at (0,2) {$2$};
			\end{tikzpicture}
		\end{center}
		Xét tính đúng hay sai của các mệnh đề sau:
		\begin{enumerate}[label=\alph*)]
			\item Hàm số đã cho đồng biến trên khoảng $(-1; 1)$.
			\item Đồ thị hàm số nhận điểm giao của đồ thị với trục tung là $I(0; 2)$ làm tâm đối xứng.
			\item Hàm số đạt cực đại tại $x = 4$ và đạt cực tiểu tại $x = 0$.
			\item Đường thẳng đi qua hai điểm cực trị của đồ thị hàm số có phương trình $y = 2x + 2$.
		\end{enumerate}
		
		\vspace{4mm}
		% CÂU 2: BÀI TOÁN THỰC TẾ (Ứng dụng đơn điệu, cực trị vào động học)
		\item Một vật chuyển động thẳng có phương trình quãng đường xác định bởi:
		\[ s(t) = -t^3 + 9t^2 + 120t \quad (t \ge 0) \]
		trong đó thời gian $t$ tính bằng giây (s) và quãng đường $s(t)$ tính bằng mét (m). Vận tốc tức thời của vật được tính bởi công thức $v(t) = s'(t)$ và gia tốc tức thời được tính bởi công thức $a(t) = v'(t)$. Xét tính đúng hay sai của các mệnh đề sau:
		\begin{enumerate}[label=\alph*)]
			\item Kể từ lúc bắt đầu chuyển động, vật di chuyển cùng chiều dương của trục tọa độ và dừng lại tại thời điểm $t = 10$ giây.
			\item Trong khoảng thời gian từ $0$ giây đến $10$ giây, vận tốc của vật luôn luôn tăng.
			\item Gia tốc tức thời của vật tại thời điểm $t = 2$ giây có giá trị bằng $6 \text{ m/s}^2$.
			\item Vận tốc của vật đạt giá trị cực đại (lớn nhất) tại thời điểm $t = 3$ giây.
		\end{enumerate}
		
	\end{enumerate}
	
	% ================================================
	% ĐỀ SỐ 2 - PHẦN III: TRẢ LỜI NGẮN
	% ================================================
	\vspace{4mm}
	\noindent\textbf{PHẦN III. Câu trắc nghiệm trả lời ngắn.} Thí sinh trả lời từ câu 1 đến câu 4.
	\vspace{2mm}
	
	\begin{enumerate}[label=\textbf{Câu \arabic*.}, leftmargin=*]
		
		% CÂU 1: ĐƠN ĐIỆU HÀM BẬC 3 (Hỏi độ dài khoảng nghịch biến)
		\item Biết hàm số $y = \dfrac{1}{3}x^3 - x^2 - 3x + 10$ nghịch biến trên khoảng $(a; b)$. Tính giá trị của biểu thức $T = b - a$.
		
		% CÂU 2: CỰC TRỊ HÀM BẬC 3 (Tích hai giá trị cực trị)
		\item Cho hàm số $y = -x^3 + 3x^2 + 1$. Tính tích của giá trị cực đại và giá trị cực tiểu của hàm số đã cho.
		
		% CÂU 3: BÀI TOÁN THỰC TẾ (Nồng độ thuốc trong máu - Cực đại)
		\item Sau khi uống một viên thuốc, nồng độ thuốc trong máu của một bệnh nhân sau $t$ giờ được mô hình hóa bởi hàm số $C(t) = -t^3 + 9t^2 + 21t$ (với $t \ge 0$, nồng độ $C(t)$ tính bằng mg/l). Hỏi sau khi uống thuốc bao nhiêu giờ thì nồng độ thuốc trong máu của bệnh nhân đạt giá trị lớn nhất?
		
		% CÂU 4: MAX MIN HÀM BẬC 2/BẬC NHẤT TRÊN ĐOẠN
		\item Tìm giá trị nhỏ nhất của hàm số $y = \dfrac{x^2 - x + 4}{x - 1}$ trên đoạn $[2; 4]$.
		
	\end{enumerate}
	
	% ================================================
	% ĐỀ SỐ 2 - PHẦN IV: TỰ LUẬN
	% ================================================
	\vspace{4mm}
	\noindent\textbf{PHẦN IV. Tự luận.} Thí sinh trình bày bài giải chi tiết từ bài 1 đến bài 3.
	\vspace{2mm}
	
	\begin{enumerate}[label=\textbf{Bài \arabic*.}, leftmargin=*]
		
		% BÀI 1: BÀI TOÁN THỰC TẾ (Tìm điểm cực tiểu của hàm chi phí vận hành)
		\item Một tàu chở hàng chạy từ cảng A đến cảng B với vận tốc $v$ (km/h) ($v > 0$). Chi phí nhiên liệu cho chuyến đi (tính bằng triệu đồng) được mô hình hóa bởi hàm số:
		\[ F(v) = 0.2v^3 - 9v^2 + 120v + 1000 \]
		Tìm vận tốc $v$ để chi phí nhiên liệu của chuyến hàng đạt cực tiểu (giá trị nhỏ nhất cục bộ của hàm chi phí).
		
		\vspace{2mm}
		% BÀI 2: CỰC TRỊ NÂNG CAO (Tính diện tích tam giác cực trị với gốc tọa độ O)
		\item Cho hàm số $y = \dfrac{x^2 - 3x + 6}{x - 1}$ có đồ thị là đường cong $(C)$. Gọi $A, B$ lần lượt là hai điểm cực trị của đồ thị $(C)$. Chứng minh rằng ba điểm $O(0;0), A, B$ không thẳng hàng và tính diện tích tam giác $OAB$.
		
		\vspace{2mm}
		% BÀI 3: KHẢO SÁT HÀM SỐ (Khảo sát hàm đa thức bậc ba đầy đủ)
		\item Khảo sát sự biến thiên (không vẽ bảng biến thiên), đồng thời xác định tâm đối xứng của hàm số $y = -x^3 + 3x^2 - 4$
		
	\end{enumerate}
	
	
	% ================================================
	% ĐỀ SỐ 3
	% ================================================
	\newpage % Bắt buộc sang trang mới
	\begin{center}
		{\LARGE \textbf{ĐỀ SỐ 3}}\\[4pt]
	\end{center}
	\vspace{4mm}
	
	% ================================================
	% ĐỀ SỐ 3 - PHẦN I: TRẮC NGHIỆM
	% ================================================
	\noindent\textbf{PHẦN I. Câu trắc nghiệm nhiều phương án lựa chọn.} Thí sinh trả lời từ câu 1 đến câu 12. Mỗi câu hỏi thí sinh chỉ chọn một phương án.
	\vspace{2mm}
	
	\begin{enumerate}[label=\textbf{Câu \arabic*.}, leftmargin=*]
		
		% CÂU 1: ĐẠO HÀM (Hàm phân thức bậc 1/bậc 1 đổi hệ số)
		\item Cho hàm số $y = \dfrac{3 - x}{2x + 1}$. Đạo hàm $y'$ của hàm số đã cho là:
		\begin{multicols}{2}
			\begin{enumerate}[label=\Alph*.]
				\item $y' = \dfrac{5}{(2x+1)^2}$.
				\item $y' = \dfrac{-5}{(2x+1)^2}$.
				\item $y' = \dfrac{-7}{(2x+1)^2}$.
				\item $y' = \dfrac{7}{(2x+1)^2}$.
			\end{enumerate}
		\end{multicols}
		
		% CÂU 2: ĐƠN ĐIỆU - Nhận biết từ bảng biến thiên ngược
		\item Cho hàm số $y = f(x)$ liên tục trên $\mathbb{R}$ và có bảng biến thiên dưới đây:
		
		\begin{center}
			\begin{tikzpicture}[xscale=1.1, yscale=0.8]
				\draw (0,0) rectangle (8,3);
				\draw (0,1) -- (8,1);
				\draw (0,2) -- (8,2);
				\draw (1,0) -- (1,3);
				
				\node at (0.5, 2.5) {$x$};
				\node at (1.5, 2.5) {$-\infty$};
				\node at (3.5, 2.5) {$0$};
				\node at (5.5, 2.5) {$4$};
				\node at (7.5, 2.5) {$+\infty$};
				
				\node at (0.5, 1.5) {$y'$};
				\node at (2.5, 1.5) {$-$};
				\node at (3.5, 1.5) {$0$};
				\node at (4.5, 1.5) {$+$};
				\node at (5.5, 1.5) {$0$};
				\node at (6.5, 1.5) {$-$};
				
				\node at (0.5, 0.5) {$y$};
				\node at (1.5, 0.7) {$+\infty$};
				\node at (3.5, 0.2) {$-3$};
				\node at (5.5, 0.7) {$5$};
				\node at (7.5, 0.2) {$-\infty$};
				
				\draw[->, thick, blue] (1.8, 0.6) -- (3.2, 0.3);
				\draw[->, thick, blue] (3.8, 0.3) -- (5.2, 0.6);
				\draw[->, thick, blue] (5.8, 0.6) -- (7.2, 0.3);
			\end{tikzpicture}
		\end{center}
		
		Hàm số đã cho đồng biến trên khoảng nào dưới đây?
		\begin{center}
			\begin{tabular}{p{0.47\textwidth} p{0.47\textwidth}}
				A. $(0;4)$. &
				B. $(-\infty;0)$ và $(4;+\infty)$. \\[2mm]
				C. $(-3;5)$. &
				D. $(0;5)$.
			\end{tabular}
		\end{center}
		
		% CÂU 3: ĐƠN ĐIỆU - Thông hiểu hàm bậc 2 / bậc nhất
		\item Hàm số $y = \dfrac{x^2 - x + 4}{x - 1}$ đồng biến trên các khoảng nào dưới đây?
		\begin{enumerate}[label=\Alph*.]
			\item $(-1; 3)$.
			\item $(-\infty; -1) \cup (3; +\infty)$.
			\item $(-1; 1)$ và $(1; 3)$.
			\item $(-\infty; -1)$ và $(3; +\infty)$.
		\end{enumerate}
		
		% CÂU 4: ĐƠN ĐIỆU - Vận dụng hàm đa thức bậc 3 hệ số a âm
		\item Tìm các khoảng nghịch biến của hàm số bậc ba $y = -\dfrac{1}{3}x^3 + x^2 + 8x - 5$.
		\begin{enumerate}[label=\Alph*.]
			\item $(-2; 4)$.
			\item $(-\infty; -2)$ và $(4; +\infty)$.
			\item $(-\infty; -2) \cup (4; +\infty)$.
			\item $(-\infty; 4)$ và $(-2; +\infty)$.
		\end{enumerate}
		
		% CÂU 5: CỰC TRỊ - Nhận biết từ đồ thị hàm số đi xuống đi lên
		\item Cho hàm số bậc ba $y = f(x)$ liên tục trên $\mathbb{R}$ và có đồ thị như hình vẽ bên dưới.
		\begin{center}
			\begin{tikzpicture}[scale=0.8]
				\draw[->] (-2,0) -- (3,0) node[right] {$x$};
				\draw[->] (0,-3) -- (0,3) node[above] {$y$};
				\node[below, left] at (0,0) {$O$};
				% Vẽ hàm bậc 3 y = x^3 - 1.5x^2 - 1 có cực trị tại x=0, x=1
				\draw[domain=-1.1:2.1, smooth, variable=\x, blue, very thick] plot ({\x}, {\x*\x*\x - 1.5*\x*\x - 1});
			\end{tikzpicture}
		\end{center}
		Số điểm cực trị của đồ thị hàm số đã cho là:
		\begin{multicols}{4}
			\begin{enumerate}[label=\Alph*.]
				\item $2$.
				\item $1$.
				\item $3$.
				\item $0$.
			\end{enumerate}
		\end{multicols}
		
		% CÂU 6: CỰC TRỊ - Thông hiểu tọa độ điểm cực tiểu hàm bậc 2 / bậc 1
		\item Tìm tọa độ điểm cực tiểu của đồ thị hàm số $y = \dfrac{x^2 + 3}{x - 1}$.
		\begin{multicols}{4}
			\begin{enumerate}[label=\Alph*.]
				\item $(-1; -2)$.
				\item $(3; 6)$.
				\item $x = 3$.
				\item $y = 6$.
			\end{enumerate}
		\end{multicols}
		
		% CÂU 7 MỚI: CỰC TRỊ HÀM BẬC 3 (Đường thẳng đi qua 2 điểm cực trị - Mức độ Vận dụng)
		\item Biết đồ thị hàm số $y = x^3 - 3x^2 - 9x + 2$ có hai điểm cực trị $A, B$. Đường thẳng $AB$ đi qua điểm nào dưới đây?
		\begin{multicols}{4}
			\begin{enumerate}[label=\Alph*.]
				\item $M(1; -1)$.
				\item $N(0; 2)$.
				\item $P(0; -1)$.
				\item $Q(-1; -7)$.
			\end{enumerate}
		\end{multicols}
		
		% CÂU 8: MAX MIN - Nhận biết từ đồ thị tiệm cận bị chặn
		\item Gọi $M, m$ lần lượt là giá trị lớn nhất và nhỏ nhất của hàm số trên đoạn $[2; 4]$ có đồ thị (đường cong màu xanh) như hình vẽ bên dưới.
		\begin{center}
			\begin{tikzpicture}[scale=0.8]
				\draw[->] (-0.5,0) -- (5,0) node[right] {$x$};
				\draw[->] (0,-0.5) -- (0,6) node[above] {$y$};
				\node[below, left] at (0,0) {$O$};
				
				% Tiệm cận đứng nét đứt màu đỏ
				\draw[red, dashed, thick] (1,0) -- (1,5.5);
				% Nhánh đồ thị
				\draw[domain=1.4:4.5, smooth, variable=\x, blue, very thick] plot ({\x}, {(2*\x + 1)/(\x - 1)});
				
				% Gióng tọa độ
				\draw[dashed, gray] (2,0) -- (2,5) -- (0,5);
				\draw[dashed, gray] (4,0) -- (4,3) -- (0,3);
				\fill[red] (2,5) circle (2pt);
				\fill[red] (4,3) circle (2pt);
				
				\node[below, red] at (2,0) {$2$};
				\node[below, red] at (4,0) {$4$};
				\node[left, red] at (0,5) {$5$};
				\node[left, red] at (0,3) {$3$};
			\end{tikzpicture}
		\end{center}
		Tính giá trị của biểu thức $T = M + m$.
		\begin{multicols}{4}
			\begin{enumerate}[label=\Alph*.]
				\item $10$.
				\item $4$.
				\item $6$.
				\item $8$.
			\end{enumerate}
		\end{multicols}
		
		% CÂU 9: MAX MIN - Thông hiểu hàm bậc 3 trên đoạn
		\item Tìm giá trị nhỏ nhất của hàm số $y = x^3 - 3x^2 - 9x + 2$ trên đoạn $[-2; 2]$.
		\begin{multicols}{4}
			\begin{enumerate}[label=\Alph*.]
				\item $0$.
				\item $-10$.
				\item $-20$.
				\item $7$.
			\end{enumerate}
		\end{multicols}
		
		% CÂU 10: TIỆM CẬN (Thông hiểu phân tích tiệm cận hàm bậc 2/ bậc 1)
		\item Đường tiệm cận đứng của đồ thị hàm số $y = \dfrac{x^2 - 4x + 3}{x - 2}$ là đường thẳng có phương trình:
		\begin{multicols}{4}
			\begin{enumerate}[label=\Alph*.]
				\item $x = 2$.
				\item $y = x - 2$.
				\item $x = 1$.
				\item Không có tiệm cận đứng.
			\end{enumerate}
		\end{multicols}
		
		% CÂU 11: TÂM ĐỐI XỨNG CỦA HÀM SỐ (Thông hiểu hàm bậc 2 / bậc nhất)
		\item Tìm tọa độ tâm đối xứng $I$ (giao điểm của hai đường tiệm cận) của đồ thị hàm số $y = \dfrac{x^2 - x + 2}{x - 1}$.
		\begin{multicols}{4}
			\begin{enumerate}[label=\Alph*.]
				\item $I(1; 0)$.
				\item $I(0; 1)$.
				\item $I(1; 1)$.
				\item $I(-1; -1)$.
			\end{enumerate}
		\end{multicols}
		
		% CÂU 12: NHẬN DIỆN HÀM SỐ (Thông hiểu đồ thị hàm đa thức cực trị đẹp)
		\item Đường cong trong hình vẽ dưới đây là đồ thị của hàm số nào?
		\begin{center}
			\begin{tikzpicture}[scale=0.8]
				\draw[->] (-2.5,0) -- (2.5,0) node[right] {$x$};
				\draw[->] (0,-3) -- (0,3) node[above] {$y$};
				\node[below, left] at (0,0) {$O$};
				% Hàm y = -x^3 + 3x
				\draw[domain=-1.8:1.8, smooth, variable=\x, blue, very thick] plot ({\x}, {-\x*\x*\x + 3*\x});
				\draw[dashed, gray] (-1,0) -- (-1,-2) -- (0,-2);
				\draw[dashed, gray] (1,0) -- (1,2) -- (0,2);
				\node[above] at (-1,0) {$-1$};
				\node[below] at (1,0) {$1$};
				\node[left] at (0,2) {$2$};
				\node[right] at (0,-2) {$-2$};
			\end{tikzpicture}
		\end{center}
		\begin{multicols}{2}
			\begin{enumerate}[label=\Alph*.]
				\item $y = x^3 - 3x$.
				\item $y = -x^3 + 3x$.
				\item $y = -x^3 - 3x$.
				\item $y = -x^3 + 3x^2$.
			\end{enumerate}
		\end{multicols}
		
	\end{enumerate}
	
	% ================================================
	% ĐỀ SỐ 3 - PHẦN II: ĐÚNG / SAI
	% ================================================
	\vspace{4mm}
	\noindent\textbf{PHẦN II. Câu trắc nghiệm đúng sai.} Thí sinh trả lời từ câu 1 đến câu 2. Trong mỗi ý a), b), c), d) ở mỗi câu, thí sinh chọn đúng hoặc sai.
	\vspace{2mm}
	
	\begin{enumerate}[label=\textbf{Câu \arabic*.}, leftmargin=*]
		
		% CÂU 1 ĐÚNG/SAI: KHẢO SÁT ĐỒ THỊ (Hàm phân thức bậc 1/bậc 1)
		\item Cho hàm số $y = f(x) = \dfrac{2x - 1}{x + 1}$ có đồ thị như hình vẽ bên dưới.
		\begin{center}
			\begin{tikzpicture}[scale=0.8]
				\draw[->] (-4,0) -- (4,0) node[right] {$x$};
				\draw[->] (0,-4) -- (0,5) node[above] {$y$};
				\node[below, left] at (0,0) {$O$};
				
				% Tiệm cận đứng và tiệm cận ngang dạng nét đứt màu đỏ
				\draw[red, dashed, thick] (-1,-4) -- (-1,5);
				\draw[red, dashed, thick] (-4,2) -- (4,2);
				\node[below, red] at (-1,0) {$-1$};
				\node[left, red] at (0,2) {$2$};
				
				% Hai nhánh đồ thị hàm số y = (2x-1)/(x+1)
				\draw[domain=-4:-1.45, smooth, variable=\x, blue, very thick] plot ({\x}, {(2*\x - 1)/(\x + 1)});
				\draw[domain=-0.65:4, smooth, variable=\x, blue, very thick] plot ({\x}, {(2*\x - 1)/(\x + 1)});
				
				% Giao điểm đặc biệt với Ox, Oy
				\fill (0,-1) circle (1.5pt) node[right] {$-1$};
				\fill (0.5,0) circle (1.5pt) node[below] {$0.5$};
			\end{tikzpicture}
		\end{center}
		Xét tính đúng hay sai của các mệnh đề sau:
		\begin{enumerate}[label=\alph*)]
			\item Hàm số luôn đồng biến trên tập xác định $\mathbb{R} \setminus \{-1\}$.
			\item Đồ thị hàm số nhận giao điểm của hai đường tiệm cận là điểm $I(-1; 2)$ làm tâm đối xứng.
			\item Giao điểm của đồ thị hàm số với trục hoành có hoành độ bằng $0{,}5$.
			\item Đường tiệm cận ngang của đồ thị hàm số là đường thẳng $y = 2$.
		\end{enumerate}
		
		\vspace{4mm}
		% CÂU 2 ĐÚNG/SAI: BÀI TOÁN THỰC TẾ (Hóa học - Tốc độ phản ứng)
		\item Nồng độ của một sản phẩm hóa học $C(t)$ (đơn vị: mg/l) sau thời gian $t$ phút kể từ lúc phản ứng bắt đầu ($0 \le t \le 15$) được mô hình hóa bởi hàm số bậc ba:
		\[ C(t) = -\dfrac{1}{3}t^3 + 4t^2 + 9t \]
		Tốc độ thay đổi nồng độ sản phẩm tại thời điểm $t$ là đạo hàm bậc nhất của nồng độ theo thời gian: $v(t) = C'(t)$. Xét tính đúng hay sai của các mệnh đề sau:
		\begin{enumerate}[label=\alph*)]
			\item Nồng độ sản phẩm $C(t)$ đạt giá trị cực đại (lớn nhất) tại thời điểm $t = 9$ phút kể từ lúc phản ứng bắt đầu.
			\item Tốc độ thay đổi nồng độ sản phẩm $v(t)$ luôn luôn giảm trong suốt quá trình phản ứng từ thời điểm $0$ phút đến $15$ phút.
			\item Tốc độ thay đổi nồng độ sản phẩm đạt giá trị cực đại tại thời điểm $t = 4$ phút.
			\item Tại thời điểm $t = 2$ phút, tốc độ thay đổi nồng độ sản phẩm đang có xu hướng tăng dần.
		\end{enumerate}
		
	\end{enumerate}
	% ================================================
	% ĐỀ SỐ 3 - PHẦN III: TRẢ LỜI NGẮN
	% ================================================
	\vspace{4mm}
	\noindent\textbf{PHẦN III. Câu trắc nghiệm trả lời ngắn.} Thí sinh trả lời từ câu 1 đến câu 4.
	\vspace{2mm}
	
	\begin{enumerate}[label=\textbf{Câu \arabic*.}, leftmargin=*]
		
		% CÂU 1: ĐƠN ĐIỆU HÀM BẬC 3 (Tổng các số nguyên trong khoảng đơn điệu)
		\item Biết hàm số $y = \dfrac{1}{3}x^3 - x^2 - 3x + 10$ nghịch biến trên khoảng $(a; b)$. Tính giá trị của biểu thức $T = b - a$.
		
		% CÂU 2: CỰC TRỊ HÀM BẬC 2 / BẬC 1 (Tính tích cực đại, cực tiểu)
		\item Cho hàm số $y = \dfrac{x^2 - x + 4}{x - 1}$. Tính tổng các giá trị cực đại và cực tiểu của hàm số đã cho.
		
		% CÂU 3: BÀI TOÁN THỰC TẾ (Nồng độ thuốc trong máu - Cực đại)
		\item Sau khi uống một viên thuốc, nồng độ thuốc trong máu của một bệnh nhân sau $t$ giờ được mô hình hóa bởi hàm số $C(t) = -t^3 + 9t^2 + 21t$ (với $t \ge 0$, nồng độ $C(t)$ tính bằng mg/l). Hỏi sau khi uống thuốc bao nhiêu giờ thì nồng độ thuốc trong máu của bệnh nhân đạt giá trị lớn nhất?
		
		% CÂU 4: MAX MIN HÀM BẬC 1/BẬC 1 TRÊN ĐOẠN (Tính tổng Max, Min)
		\item Tìm tổng giá trị lớn nhất và giá trị nhỏ nhất của hàm số $y = \dfrac{x - 3}{2x - 1}$ trên đoạn $[1; 3]$.
		
	\end{enumerate}
	
	% ================================================
	% ĐỀ SỐ 3 - PHẦN IV: TỰ LUẬN
	% ================================================
	\vspace{4mm}
	\noindent\textbf{PHẦN IV. Tự luận.} Thí sinh trình bày bài giải chi tiết từ bài 1 đến bài 3.
	\vspace{2mm}
	
	\begin{enumerate}[label=\textbf{Bài \arabic*.}, leftmargin=*]
		
		% BÀI 1: BÀI TOÁN THỰC TẾ (Công suất tối đa của tuabin khí)
		\item Một tuabin phát điện có hiệu suất sinh năng lượng $P$ (kW) phụ thuộc vào tốc độ quay $v$ (vòng/phút) ($v \ge 0$) theo mô hình hàm bậc ba:
		\[ P(v) = -0.1v^3 + 4.5v^2 \]
		Xác định tốc độ quay $v$ của tuabin để công suất năng lượng phát ra đạt giá trị lớn nhất.
		
		\vspace{2mm}
		% BÀI 2: CỰC TRỊ NÂNG CAO (Khoảng cách từ O đến đường thẳng đi qua cực trị bậc 3)
		\item Cho hàm số $y = x^3 - 3x^2 - 9x + 5$ có đồ thị là đường cong $(C)$. Biết đồ thị $(C)$ có hai điểm cực trị là $A$ và $B$. Viết phương trình đường thẳng đi qua hai cực trị $A, B$ và tính khoảng cách từ gốc tọa độ $O(0;0)$ đến đường thẳng đó.
		
		\vspace{2mm}
		% BÀI 3: KHẢO SÁT HÀM SỐ (Khảo sát hàm bậc nhất / bậc nhất đầy đủ)
		\item Khảo sát sự biến thiên (không vẽ bảng biến thiên), đồng thời xác định tâm đối xứng của hàm số $y = \dfrac{2x - 3}{x - 1}$
		
	\end{enumerate}
	
\end{document}